%% SPLASH/ISSTA 2026 Doctoral Symposium Submission
%% ACM SIGPLAN acmart style, double-column, max 3 pages incl. references

\documentclass[sigplan,10pt,anonymous=false]{acmart}
\settopmatter{printfolios=false,printccs=false,printacmref=false}

\renewcommand\footnotetextcopyrightpermission[1]{}
\pagestyle{plain}

\usepackage{xspace}
\usepackage{enumitem}
\setlist{noitemsep, topsep=3pt}

\begin{document}

\title{Toward Reliable Automatic Proxies for Code Comprehensibility}

\author{Erfan Arvan}
\affiliation{%
  \institution{New Jersey Institute of Technology}
  \city{Newark}
  \state{NJ}
  \country{USA}
}
\email{ea442@njit.edu}

\begin{abstract}
Improving code comprehensibility requires measuring it reliably, yet no
reliable automatic proxy for comprehensibility exists.
Developing one requires a validated human-study baseline to test against.
However, the human-study proxies used in the literature (e.g., Likert-scale
ratings or response time on output-prediction tasks) have not been
systematically validated either, and the construct itself lacks a standard
definition.
We address this foundational gap first: we introduce a conceptualization of
comprehensibility based on expert agreement, construct a ground truth via the
Delphi protocol, and evaluate commonly used human-study proxies against it.
We find that proxies based on input--output reasoning and response time are
particularly reliable, while syntax-based proxies are especially unreliable.
With this validated foundation, we propose a controlled human study to test
whether code verifiability---the ease with which static verification tools
can analyze code---can serve as a reliable automatic proxy for comprehensibility.
If verifiability alone explains comprehension difficulty well, it can be used
as a standalone automatic proxy; if it only partially explains it, it becomes
one signal among several in a combined automatic proxy.
Beyond this study, the research agenda includes reassessing prior work in
light of proxy reliability, evaluating LLM-based agents as substitutes for
human participants, and developing tools that leverage reliable automatic
proxies to improve code comprehensibility.
\end{abstract}

\maketitle

%% -----------------------------------------------------------------------
\section{Motivation}
%% -----------------------------------------------------------------------

Code comprehensibility is a fundamental property of software: most
maintenance, debugging, refactoring, and code review tasks depend on a
developer's ability to build an accurate mental model of the code they are
working with.
Developers spend between 58\% and 70\% of their time on this
activity~\cite{percent1,scalabrino2017automatically}, and code that is
difficult to understand slows development, increases the risk of introducing
bugs, and makes collaboration harder.
Improving code comprehensibility is therefore a high-impact goal---but
improving it requires measuring it reliably.

Most existing studies measure comprehensibility through controlled human
studies, which require careful design, participant recruitment, and often
months of preparation.
This makes it infeasible to rely on human studies whenever we need a
comprehensibility estimate---for instance, to compare two implementations,
to guide automated refactoring, or to evaluate LLM-generated code for
human understandability.
What we need instead is a reliable \emph{automatic} proxy for code
comprehensibility: a program analysis or metric that estimates comprehension
difficulty directly from the code, without requiring human participants.
Prior proposals---cyclomatic complexity~\cite{curtis2006measuring}, Halstead measures~\cite{chidamber1994metrics}, lines of
code~\cite{nunez2017source}---turn out to correlate weakly, if at all, with developer comprehension
behavior~\cite{feigenspan2011exploring,peitek2021program,scalabrino2019automatically}.
No reliable automatic proxy for code comprehensibility currently exists.

These prior proposals are surface-level and syntax-based: they capture
structural properties of code without reasoning about its meaning or behavior.
A natural next step is to ask whether semantic information can do better.
One promising direction is code \emph{verifiability}.
The hypothesis is that code that is easier for verification tools to analyze
is also easier for humans to understand.
Prior work~\cite{FeldmanKC2023} investigated this hypothesis in a
meta-analysis, operationalizing verifiability as the number of false positive
warnings produced by static analysis tools, and found associations between
verifiability and comprehension measures across a collection of prior studies.
However, because their analysis aggregated heterogeneous prior studies with
different code sources, structures, and participant populations, they could
not control for confounders or establish a causal relationship.
We want to re-evaluate this relationship in a controlled setting where
verifiability is the only factor that varies between code snippets.

To do this rigorously, we need a reliable human-study baseline to test
verifiability against: we need to know which code is actually harder to
understand from a human study.
This is where we encountered a first obstacle.
The human-study proxies used in the literature---observable tasks assigned
to participants (e.g., predicting program output) paired with measurements
of their performance (e.g., response time or correctness)---have themselves
never been systematically validated.
Moreover, the field lacks a widely accepted definition of comprehensibility,
so researchers often define it in terms of whatever proxy they happen to
use~\cite{Wyrich:ACS23}.
Without knowing which human-study proxies are reliable, we cannot use them
confidently to evaluate a candidate automatic proxy.

We therefore addressed this obstacle first.
We introduce a principled conceptualization of comprehensibility, construct
an expert-consensus ground truth, and evaluate the reliability of commonly
used human-study proxies against it.
With that validated foundation in hand, we then proceed to test verifiability
as a candidate automatic proxy in a controlled human study.

%% -----------------------------------------------------------------------
\section{Phase 1 (Completed): Reconceptualizing Comprehensibility and\\
Evaluating Proxy Reliability}
%% -----------------------------------------------------------------------

We addressed the first obstacle---the lack of validated human-study
proxies---in our recent work~\cite{arvan2026reliability}.
Our goal was to compare the reliability of commonly used human-study proxies
against a ground truth measure of comprehension difficulty.
The major challenge was that no such ground truth existed: there was no
independent, principled measure of which code is actually harder to understand
that we could use to evaluate the proxies against.

To construct one, we conceptualized code comprehensibility as expert agreement
on relative comprehension difficulty.
Using the Delphi protocol~\cite{delphi,delphiOriginal,randManual}---a
structured method for eliciting and refining expert judgment under
uncertainty, successfully applied in medicine~\cite{m1,m2,m3,m4},
geopolitical forecasting~\cite{g1,g2}, and other domains~\cite{e1,o1}---a
panel of five senior professional software engineers iteratively ranked eight
real-world Java methods by comprehension difficulty and resolved disagreements
through anonymized written feedback over three rounds.

We then conducted a controlled study in which 44 undergraduate CS students
completed tasks corresponding to 14 comprehension proxies from the literature
on the same eight methods, covering four task types (program function,
input--output reasoning, syntax identification, and subjective rating) and
three measure types (response time, answer correctness, and Likert-scale
ratings).
Correlating each proxy against the expert-consensus ranking via Kendall's
$\tau$ revealed a clear hierarchy:
\begin{itemize}
  \item \textbf{Most reliable:} time-based input--output proxies:\\
        (\emph{time\_output}, $\tau = 0.86$; \emph{time\_correct\_output},
        $\tau = 0.79$; \emph{time\_function}, $\tau = 0.71$).
  \item \textbf{Moderately reliable:} subjective difficulty ratings and
        general reading time ($\tau \in [0.50, 0.64]$).
  \item \textbf{Unreliable:} all syntax-based proxies
        (\emph{acc\_syntaxBL}, $\tau = -0.37$;
        \emph{time\_syntaxBL}, $\tau = -0.21$), showing near-zero or
        negative correlations regardless of whether time or correctness
        is used as the measure.
\end{itemize}

These results have two consequences.
First, they provide a validated set of proxies for future human studies:
researchers can now confidently use input--output response time as a reliable
approximation of comprehension difficulty.
Second, they call into question prior work that relied on syntax-based
proxies, whose negative correlation with expert judgments suggests they
measure a fundamentally different construct than comprehension difficulty.

%% -----------------------------------------------------------------------
\section{Phase 2 (Proposed): Verifiability as an Automatic Proxy}
%% -----------------------------------------------------------------------

With a validated human-study baseline in hand, we can now rigorously test
candidate automatic proxies.
The first candidate we investigate is code \emph{verifiability}: the ease
with which static verification tools can analyze code.
This is motivated by a widely held assumption in the verification
community---captured explicitly in the Checker Framework
manual~\cite{checkerframework-manual}, which advises developers to ``rewrite
your code to be simpler for the checker to analyze; this is likely to make
it easier for people to understand, too''---that more verifiable code is also
more comprehensible.

A prior meta-analysis~\cite{FeldmanKC2023} found associations between
false positive warnings produced by static analysis tools and comprehension measures.
However, because it aggregated heterogeneous prior studies, it cannot isolate
verifiability as the cause: the snippets differ in source, structure, size,
and---as Phase~1 shows---proxy quality.
Our proposed study addresses all of these limitations.

\paragraph{Treatment.}
Verifiability is operationalized as a binary condition relative to a fixed
set of Checker Framework checkers~\cite{PapiACPE2008}.
\emph{High-verifiability} code produces no false positive warnings;
\emph{low-verifiability} code is a semantically equivalent version of the
same method that produces one or more false positives under the same
checkers.
Paired variants are constructed via semantics-preserving transformations
(defined as input--output equivalence) in two phases: first by eliminating
false positives from warning-producing methods, then by injecting false
positives into warning-free methods.
Transformations span a range of code-length changes to avoid collinearity
between verifiability condition and code length.

\paragraph{Outcome and design.}
The primary outcome is response time on input--output tasks---the most
reliable proxy identified in Phase~1.
The experiment uses a mixed design: within-subjects over snippets and
between-subjects over verifiability condition, so no participant sees both
versions of the same method.
Latin square counterbalancing ensures balanced exposure to both conditions
across participants.
The analysis uses linear mixed-effects models with a fixed effect for
verifiability condition, random intercepts for participants and snippets,
and covariates for snippet order, code length, and transformation type.

\paragraph{Interpretation.}
If verifiability alone explains comprehension difficulty well, it can be used
as a standalone automatic proxy.
If it only partially explains comprehension difficulty, it becomes one signal
among several in a combined automatic proxy.
A null result is equally meaningful: it would provide controlled evidence
that the associations in prior work~\cite{FeldmanKC2023} were driven by
confounds rather than a genuine relationship.

%% -----------------------------------------------------------------------
\section{Evaluation and Future Directions}
%% -----------------------------------------------------------------------

\paragraph{Phase 2 evaluation.}
The primary evaluation metric is the fixed effect of verifiability condition on
response time on input--output tasks.
The main threats to validity are:
\textit{internal}---transformations may introduce residual structural
differences beyond verifiability, mitigated by simple, manually validated
transformations and structural covariates in the model;
\textit{construct}---false positive warnings and response time are each
one-dimensional operationalizations of broader constructs; and
\textit{external}---results are scoped to Java, the Checker Framework, and
undergraduate participants, and generalization to other languages, tools,
and professional developers is future work.
All data and analysis scripts will be made publicly available.

\paragraph{Broader agenda.}
Beyond Phase~2, the research agenda includes:
(1)~evaluating whether LLM-based agents can reliably substitute for human
participants in comprehension studies, enabling further scaling of proxy
validation;
(2)~systematically reassessing prior comprehensibility studies through the
lens of validated proxies, identifying which findings remain robust and which
relied on unreliable measurements;
(3)~investigating how different forms of accidental
complexity~\cite{brooks1987no} affect comprehension time, using input--output
response time as the validated outcome; and
(4)~developing tools that identify hard-to-understand code and suggest
semantics-preserving refactorings, contingent on establishing a reliable
automatic proxy.

Together, this work aims to establish a principled and empirically grounded
foundation for measuring and improving code comprehensibility.

%% -----------------------------------------------------------------------

\bibliographystyle{ACM-Reference-Format}
\bibliography{local,delphi,plume-bib/ernst,plume-bib/crossrefs}

\end{document}